\chapter{Where One Knob Isn't Enough}

\epi{A book about counting knobs that doesn't count its own\\ is a sales brochure.}

\fig[1.0]{plate_wall}{The wall. Each red brick is a place where one knob isn't enough, and every one of them is in this chapter: clusters, the microwave background, a stubborn dwarf, the Milky Way's vertical pull, stellar masses, disc weighings, the Solar System without a screen, the matter term, and a published high-redshift rise.}


\section{Reading the row to the right}

Here is the programme's final scorecard, FG1, frozen on 1 November 2026. It grades five candidates, four versions of a cold sector the programme built and then retired, and the law alone (C4), on eleven lines. I'm showing all five because the cold sector was my attempt to cover exactly the places where the law fails, and it failed too.

\fig{fg1_heatmap}{The final grade. Row C4 is the law with nothing added. Rows C1--C3 and C5 are versions of a cold dark energy sector the programme built to cover clusters and the CMB, graded on measured data and retired on 1 November 2026. No row passes everything. The lines aren't independent, so no winner is declared.\label{fig:fg1}}

The law's row passes the galaxy lines and fails six: the cluster keys, X-COP, the Local Group dwarfs, the Milky Way, the DiskMass discs and IC~1613. Add the CMB, which isn't on this table because the law can't even compute it, and the Solar System, where the law needs a screen it doesn't contain. Let's go through them.

\section{Loss 1: galaxy clusters}

Clusters are the oldest problem for modified dynamics, older than this programme by decades, and the law inherits it whole.

\textbf{X-COP.} For seven clusters with excellent X-ray data, the hydrostatic mass, the mass needed to hold the hot gas up against its own pressure, is \textbf{2.93 times} what the law provides from the visible baryons (alternative footing 2.72). All seven come up short. If the hydrostatic masses are biased high by $0.12\dex$, a common worry, the shortfall is still 2.23 (2.06).

\textbf{Shear around 318 clusters.} Using measured gas from eROSITA and raw lensing shear from KiDS, the law is short by a factor $A=1.91\pm0.12$ between 0.5 and 1.2\,Mpc (alternative $1.75\pm0.11$) and matches, $A=1.01\pm0.09$, between 1.2 and 3\,Mpc. On the measured chain it fails 10 of 15 keys on both footings.

\fig{clusters}{Clusters. Left: on X-COP, hydrostatic masses exceed the law's by a factor near 3, and near 2 even if the hydrostatic masses are biased. Right: lensing shear around 318 eROSITA clusters needs about twice the law inside 1.2\,Mpc and agrees further out.}

\textbf{What's missing, described.} Inside about $R_{500}$ the deficit is near 8 times the gas mass at 0.2--0.5\,$R_{500}$ and 5 times at 0.5--1\,$R_{500}$, with no detected mass dependence and no preferred scale found. The outer shear grows with gas mass as $M_{\rm gas}^{0.90\pm0.21}$, against the law's 0.47 and Newton's 0.84. So the agreement in the outer bin is an average over a slope the law gets wrong.

\fig{cluster_deficit}{What clusters are missing, relative to their hot gas. Inside about $R_{500}$ the law leaves a deficit near eight times the gas mass in the inner bin and five times further out. That's a description, not a mechanism.}



\fig{outer_shear_slope}{How outer shear grows with gas mass. The measured slope is $0.90\pm0.21$; the law predicts 0.47, Newton 0.84. The law's good average amplitude at large radius hides a slope it gets wrong.}



\begin{knobbox}{A knob I tried, and didn't adopt}
The obvious patch is a gas term: boost the law in proportion to the local gas acceleration,
\[
g=g_{\rm law}\left(1+\theta\,k\,\frac{g_{\rm gas}}{a_0}\right),
\]
with $k$ fitted on clusters. With $k=13.29$ (alternative 14.25), fitted on four X-COP clusters, the canonical cluster keys go from 10 of 15 to 14. But that's fragile: it holds only for $k$ between 12.87 and 15.09, no value gives 15 of 15, it breaks discs unless a switch turns it off there (and then it passes largely by construction), and there's no action behind it. On the Bullet Cluster's lensing map, every version tied to present gas puts the mass in the wrong place, because the observed lensing peaks sit on the galaxies, not on the gas. The frozen label is KEYS FAIL. That would have been knob number two, and it didn't earn its place.
\end{knobbox}

\fig{k_window}{Why the gas-term knob failed. Each bar is the range of $k$ that one cluster test allows. One test needs $k$ below 7.2; the others need it above 8.8 or 12.9. No single setting satisfies all three, so no value of the knob gives 15 of 15.}



\section{Loss 2: the cosmic microwave background}

This is the big one. The law's scale $a_0$ doesn't appear in the linear equations of a uniform expanding universe, so the law can't change the background expansion. It can act on the perturbations, and the potentials that CMB lensing maps between redshifts 0.5 and 2.33 have gradients of 0.003 to 0.22\,$a_0$, inside the law's regime. So the question is fair.

The answer is bad. Run on the cosmic web from a start with baryons only and no cold component, the law gets the CMB lensing spectrum wrong in \emph{shape}, not just size. As QUMOND on the peculiar field it gives about 14 times the measured lensing at multipole $L\approx28$ and 0.08 times at $L\approx377$. A baryons-only start has had its small-scale fluctuations erased by photon diffusion (Silk damping), and the law boosts the largest scales most. Letting $a_0$ track evolving dark energy, as the law says it should, lowers the whole spectrum by about half without changing its shape. A screening term put in by hand removes the large-scale excess but leaves the spectrum 1.5 to 7 times low above $L\approx100$.

\fig{cmb_lensing_shape}{The law alone against CMB lensing: too much on large scales, far too little on small ones. Two points from the graded run, joined to show the shape. Tracking dark energy halves it (gold); hand-added screening (blue band) still falls short above $L\approx100$. ``1'' would be a match.}

\begin{honestbox}[The record has no CMB pass without a dust-like component]
With a component that behaves at linear order like cold dark matter ($w=0$, no sound speed), a window exists. The programme built one, a gravitating ``cold dark energy'' sector, explicitly not a particle, and graded four versions on measured data. All four fail somewhere (Figure~\ref{fig:fg1}). Two walls stopped it: whatever behaves as dust in the CMB falls into spiral galaxies and contracts with them, which the rotation curves don't allow; and the outer shear around clusters, on measured gas, misses at $Z=+4.41$. The sector was retired on 1 November 2026. \textbf{The law offers no replacement, and I don't claim it needs none.} This is exactly where \LCDM's six knobs shine.
\end{honestbox}

\section{Loss 3: lensing around galaxies, in detail}

Chapter~11 counted the \emph{shape} of galaxy lensing as a win. The \emph{amplitude} isn't. Inside 0.3\,Mpc the data sit above the law by $+0.0887\dex$ (alternative $+0.0500$) under every grader on record, and no single swap of estimator brings the full-profile grade inside $|Z|=2$. The hot gas measured through the kinematic Sunyaev--Zel'dovich effect, mostly at 200--500\,kpc, doesn't close the gap.

The amplitude also depends on colour. Red lenses read well above the law and blue lenses slightly below. One gas rule for both colours can't bring red within 2$\sigma$ without pushing blue over.

\fig{lens_colour}{Galaxy lensing amplitude relative to the law, for red and blue lenses, as more of the baryon inventory is counted. Red stays 4.87$\sigma$ above 1 even with measured hot gas. No tested model fits these profiles in absolute terms.}

\section{Loss 4: IC 1613}

One nearby dwarf galaxy, IC~1613, doesn't fit. The law's offset is $E=-0.593$, $Z=-7.9$, and it fails under every inclination the data allow. That $Z$ has no velocity error in it; with 15 per cent on the velocities the base is $-3.94$, and $-2.35$ to $-2.83$ at the central inclinations of an independent analysis. Only a much more face-on disc would close it: $E=0$ needs an inclination of about $22^\circ$, against central estimates of 34 to 38. A disturbed gas disc is a live explanation that nobody has computed. One galaxy disfavours a law; it doesn't exclude it. But it's on the list.

\section{Loss 5: the Local Group's dwarfs}

The satellites of the Milky Way and Andromeda are the most dark-matter-dominated systems known, and several classes of them miss.

\fig{dwarf_z}{Local Group dwarf classes against the law ($Z$; inside $\pm2$ passes). Ultra-faint satellites of the Milky Way miss badly, and including the host galaxy's external field, which modified dynamics says should matter, makes it worse.}

The Milky Way's ultra-faints read $Z=+5.22$; Andromeda's bright satellites $+2.20$ and its others $+2.29$. Adding the external field from the host, which in MOND-type theories weakens a satellite's internal boost, makes it worse: $+6.76$ and $+3.97$. Worse still, at fixed luminosity the data are flat in host field ($-0.05\pm0.19$ per dex) where the external-field theory predicts $-0.76$. Listed systematics (binary stars, tides, membership, heating) don't close it: a central stack of them leaves the ultra-faints at $+5.08$.

\section{Loss 6: the Milky Way's vertical pull}

The law gives the in-plane force. To test the Milky Way, you also need the force perpendicular to the disc, $K_z$, which the law as written doesn't supply. The two standard field-equation forms (AQUAL and QUMOND) can't reach the measured $K_z$ at the Sun's position inside the priors: the lowest they get is 75.9 against an upper line of 70.8. A simpler rule passes only in a few corners of parameter space and still misses the scale length. FG1 counts it as a fail.

\section{Loss 7: measured stellar masses}

The law asks for stellar mass-to-light ratios of 0.62 and 0.65 in galaxies rotating at 60--150\kms. Stellar masses from the GSWLC catalogue give 0.277 and 0.304 for the same galaxies, a miss at $Z=+3.4$ and $+5.6$. These are estimators against a different photometric band, not a direct measurement, and offsets, distances and inclinations don't carry the difference. The law wants about twice the stellar mass the photometry suggests.

\fig{upsilon_compare}{Stellar mass-to-light ratios in two bins of rotation speed: what the law needs (red) against what stellar masses from the GSWLC catalogue give (blue). The gaps are $Z=+3.4$ and $+5.6$.}



\section{Loss 8: disc thickness}

The DiskMass survey measured vertical stellar velocity dispersions in face-on discs, which weigh the disc directly. On the law's own form, the ratio that should be near 1 reads 1.239 (alternative 1.204), outside the 0.90--1.10 line. Corrections for which stars are being measured could move it inside, but the sources' own tables span a wide range. This one is open, and graded as a fail.

\section{Loss 9: the Solar System}

As written, the law leaves a nearly constant extra pull of order $a_0$ even in strong fields, and in the Solar System that shows up as an anomalous quadrupole in the Sun's gravitational field, which radio tracking of the Cassini spacecraft at Saturn constrains. The unscreened law fails that constraint by 2.44 to 4.59 times (alternative 2.85 to 5.19). Every unscreened action class in the record fails it, by 1.69 to 5.18 times. Only v3s, with its screening length, passes.

\fig{solar_system}{The Cassini constraint. Unscreened, the law and every unscreened action fail by factors of 2 to 5. With v3s's screen the signal drops more than seven orders of magnitude below the line. The screen is the price of the Solar System, and it's paid in v3s's extra structure.}

\section{Loss 10: the matter term}

There's a quieter problem underneath everything. The law takes $\rhoDE$ from the expansion history, and that history contains a matter term several times larger than the baryons. The law doesn't explain that term. Even the redshift band of Chapter~13 uses $\rho_m(z)=\omega_m(1+z)^3$ with a cold mass taken from the CMB. With only positivity, and no CMB or Hubble-constant input, that band's own check fails. So the cleanest prediction in this book isn't entirely free of the sector the programme retired.

\section{The 96-row list}

The programme also built a harness that grades any proposed $g$ on one list of 96 rows. The rows are correlated in blocks and many are disputed, so the counts don't rank anything. They do show how much of the universe each candidate can even attempt.

\fig{rq3_counts}{The 96-row list. Green pass, red fail, gold depends on the form, grey not computable. Newton's 10 passes show how little those rows ask. For the law, 38 rows need an action or a Boltzmann calculation it doesn't have. These counts don't rank the candidates.}

\section{What a successor must do}

From the failures, the programme wrote down twelve requirements that any theory replacing the law (or completing it) has to meet:

\begin{center}\small
\begin{tabular}{@{}rp{3.1cm}p{4.0cm}@{}}
\toprule
\# & must & the law\\
\midrule
1 & keep galaxies tight & \textbf{0.0541 dex}\\
2 & stay flat with surface brightness & \textbf{$-0.016$}\\
3 & sit on $\Upsilon$ at every speed & 0.52 (alt 0.44)\\
4 & little net contraction in spirals & not needed\\
5 & fit dwarfs & \textbf{$-0.025$}; ultra-faints $+5.08$\\
6 & fit IC 1613 & $-0.593$\\
7 & fit the Milky Way & AQUAL, QUMOND infeasible\\
8 & fit isolated lenses & shape preferred; 0.09 dex short inside 0.3 Mpc\\
9 & add mass in cluster cores & $A=1.91$\\
10 & match outer cluster shear & slope 0.47 against 0.90\\
11 & act as dust in the CMB & no cold term\\
12 & have an action & met in a scope only (v3s)\\
\bottomrule
\end{tabular}
\end{center}

The collisions are the interesting part. Rows 4 and 11 pull against each other: dust in the CMB falls into galaxies. Rows 1--3 and 10 pull against each other: the amount of extra gravity needed runs from about 0.8 of a native halo's in SPARC to 1.2--1.8 in clusters. Rows 7 and 8: systems of the same mass want different things. None of this is a theorem, but every route the programme tried ran into one of these collisions.

\fig{collision_graph}{The twelve requirements as a network. Red lines join requirements that pull against each other: dust in the microwave background against little contraction in spirals; tight galaxies against outer cluster shear; the Milky Way against isolated lenses.}



\section{The honest count}

So here's the law's control panel, with the losses included.

\begin{knobbox}{The law, all in}
One knob ($\kappa$, fitted). One shape choice ($\mu=x/(1+x)$). One adopted action (v3s) in a scope, with its screening structure and with $a_0$ still put in by hand. Galaxies: tight, flat with surface brightness, dwarfs and groups passing. Clusters: short by a factor of 2 to 3. CMB: nothing. Solar System: only with the screen. Local Group dwarfs, IC~1613, the Milky Way's vertical force, measured stellar masses, disc weighings: fail or open.
\end{knobbox}

One knob explains a remarkable amount about disc galaxies. It doesn't explain the universe. \LCDM\ explains the universe's large-scale structure with six knobs and needs dozens more to make a galaxy look like the law. My law makes a galaxy look like the law with one knob and can't make a universe.

You might think that's a draw, or a loss on points. But there's one more thing the one knob does, and it's the reason this book exists. It makes a prediction that no one can tune.
